% SOURCE_AUTHORITY: sga5_fr_workpass.tex, lines 2452--2609 (LF-normalized unit).
% SOURCE_SHA256: 791F4EFFC5E02832D5D77ED03518C8156D6F07E4C8238B03545DB93D883FBB28
\subsection*{3.3}
Sejam, para $i=1,2$, $f_i:X_i\to Y_i$ um $S$-morfismo,
\[
f=f_1\times_S f_2:X\to Y,
\qquad
f'_1=f_1\times_S Y_2,
\qquad
f'_2=Y_1\times_S f_2,
\]
\[
p'_1=X_1\times_S f_2,
\qquad
p'_2=f_1\times_S X_2.
\]
\begin{equation}
\begin{tikzcd}[column sep=3.8em,row sep=2.4em]
X_1\times_S Y_2 \arrow[d,"f'_1"'] & X \arrow[l,"p'_1"'] \arrow[dl,"f" description] \arrow[d,"p'_2"] \\
Y & Y_1\times_S X_2 \arrow[l,"f'_2"]
\end{tikzcd}
\tag{3.3.0}
\end{equation}
Sejam $L_i\in\ob\Dctf(X_i)$. Existe um único isomorfismo
\begin{equation}
f_*\uRHom_S(L_1,L_2)
\xrightarrow{\sim}
\uRHom_S(f_{1!}L_1,f_{2*}L_2)
\tag{3.3.1}
\end{equation}
cujo inverso torna comutativo o quadrado
\[
\begin{tikzcd}[column sep=4.5em,row sep=2.6em]
D(f_{1!}L_1)\overset{\mathbf L}{\otimes}_S f_{2*}L_2
  \arrow[r,"3.1.1"] \arrow[d,"(1)"']
& \uRHom_S(f_{1!}L_1,f_{2*}L_2) \arrow[d] \\
f_{1*}DL_1\overset{\mathbf L}{\otimes}_S f_{2*}L_2
  \arrow[r,"(2)"']
& f_*\uRHom_S(L_1,L_2),
\end{tikzcd}
\]
onde (1) se deduz do inverso do isomorfismo de dualidade (SGA 4 XVIII 3.1.9.6)
\[
Df_{1!}L_1\xrightarrow{\sim}f_{1*}DL_1
\]
aplicando $-\overset{\mathbf L}{\otimes}_S f_{2*}L_2$, e (2) é a composição do isomorfismo de Künneth
\[
f_{1*}DL_1\overset{\mathbf L}{\otimes}_S f_{2*}L_2
\xrightarrow{\sim}
f_*(DL_1\overset{\mathbf L}{\otimes}_S L_2)
\]
e de $f_*(3.1.1)$. Aplicando $\R\Gamma(Y,-)$, deduz-se de 3.3.1 um isomorfismo
\begin{equation}
f_*:\Hom_S(L_1,L_2)
\xrightarrow{\sim}
\Hom_S(f_{1!}L_1,f_{2*}L_2).
\tag{3.3.2}
\end{equation}
Para $u\in\Hom_S(L_1,L_2)$, a correspondência $f_*u$ será chamada \emph{imagem direta} de $u$ por $f$. Quando $Y_1=Y_2=S$, escreveremos $u_*$ em vez de $f_*u$. Se $Y_1=Y_2=S$ for o espectro de um corpo separavelmente fechado, 3.3.2 se escreverá
\begin{equation}
\Hom_S(L_1,L_2)
\xrightarrow{\sim}
\Hom\bigl(\R\Gamma_c(X_1,L_1),\R\Gamma(X_2,L_2)\bigr).
\tag{3.3.3}
\end{equation}
Assim, as correspondências cohomológicas ``atuam'' sobre a cohomologia. O enunciado seguinte permitirá interpretar essa ação de várias maneiras e, em particular, estabelecer o vínculo com as definições habituais em situações como a) ou b) anteriores.

\begin{statement}{Proposição 3.4}
Com as notações de 3.3.0, tem-se um diagrama comutativo de isomorfismos em $\Dctf(Y)$:
\begin{equation}
\begin{tikzcd}[column sep=4.0em,row sep=2.8em]
& \uRHom_S(f_{1!}L_1,f_{2*}L_2)
  \arrow[dl,"(1)"'] \arrow[dr,"(2)"] \arrow[dd,"(5)" description] & \\
f'_{1*}\uRHom_S(L_1,f_{2*}L_2) \arrow[dr,"f'_{1*}(3)"'] &&
f'_{2*}\uRHom_S(f_{1!}L_1,L_2) \arrow[dl,"f'_{2*}(4)"] \\
& f_*\uRHom_S(L_1,L_2) &
\end{tikzcd}
\tag{3.4.1}
\end{equation}
onde (5) é o inverso de 3.3.1 e as flechas (1) a (4) são definidas pelos diagramas comutativos seguintes, nos quais ``c. b.'' designa uma flecha derivada de um isomorfismo de mudança de base ou de Künneth de um dos tipos considerados no n\textsuperscript{o}~1, ``dual.'' uma flecha derivada de um isomorfismo de dualidade de tipo 2.1.4 ou (SGA 4 XVIII 3.1.9.6), e ``triv.'' um isomorfismo trivial:
\[
\begin{tikzcd}[column sep=3.7em,row sep=1.7em]
\uRHom(f_{1!}L_1\overset{\mathbf L}{\otimes}_S A_{Y_2},K_{Y_1}\overset{\mathbf L}{\otimes}_S f_{2*}L_2)
  \arrow[r,"\mathrm{c.b.}"] \arrow[d,"\mathrm{c.b.}"']
& \uRHom_S(f_{1!}L_1,f_{2*}L_2) \arrow[ddd,"(1)"] \\
\uRHom(f'_{1!}(L_1\overset{\mathbf L}{\otimes}_S A_{Y_2}),K_{Y_1}\overset{\mathbf L}{\otimes}_S f_{2*}L_2) \arrow[d,"\mathrm{dual.}"']
& \\
f'_{1*}\uRHom(L_1\overset{\mathbf L}{\otimes}_S A_{Y_2},f_1'^!(K_{Y_1}\overset{\mathbf L}{\otimes}_S f_{2*}L_2)) \arrow[d,"\mathrm{c.b.}"']
& \\
f'_{1*}\uRHom(L_1\overset{\mathbf L}{\otimes}_S A_{Y_2},K_{X_1}\overset{\mathbf L}{\otimes}_S f_{2*}L_2) \arrow[r,"\mathrm{c.b.}"']
& f'_{1*}\uRHom_S(L_1,f_{2*}L_2),
\end{tikzcd}
\]
\[
\begin{tikzcd}[column sep=3.7em,row sep=1.7em]
\uRHom(f_{1!}L_1\overset{\mathbf L}{\otimes}_S A_{Y_2},K_{Y_1}\overset{\mathbf L}{\otimes}_S f_{2*}L_2)
  \arrow[r,"\mathrm{c.b.}"] \arrow[d,"\mathrm{c.b.}"']
& \uRHom_S(f_{1!}L_1,f_{2*}L_2) \arrow[ddd,"(2)"] \\
\uRHom(f_{1!}L_1\overset{\mathbf L}{\otimes}_S A_{Y_2},f'_{2*}(K_{Y_1}\overset{\mathbf L}{\otimes}_S L_2)) \arrow[d,"\mathrm{dual.}"']
& \\
f'_{2*}\uRHom(f_2'^*(f_{1!}L_1\overset{\mathbf L}{\otimes}_S A_{Y_2}),K_{Y_1}\overset{\mathbf L}{\otimes}_S L_2) \arrow[d,"\mathrm{triv.}"']
& \\
f'_{2*}\uRHom(f_{1!}L_1\overset{\mathbf L}{\otimes}_S A_{X_2},K_{Y_1}\overset{\mathbf L}{\otimes}_S L_2) \arrow[r,"\mathrm{c.b.}"']
& f'_{2*}\uRHom_S(f_{1!}L_1,L_2),
\end{tikzcd}
\]
\[
\begin{tikzcd}[column sep=3.7em,row sep=1.7em]
\uRHom(L_1\overset{\mathbf L}{\otimes}_S A_{Y_2},K_{X_1}\overset{\mathbf L}{\otimes}_S f_{2*}L_2)
  \arrow[r,"\mathrm{c.b.}"] \arrow[d,"\mathrm{c.b.}"']
& \uRHom_S(L_1,f_{2*}L_2) \arrow[ddd,"(3)"] \\
\uRHom(L_1\overset{\mathbf L}{\otimes}_S A_{Y_2},p'_{1*}(K_{X_1}\overset{\mathbf L}{\otimes}_S L_2)) \arrow[d,"\mathrm{dual.}"']
& \\
p'_{1*}\uRHom(p_1'^*(L_1\overset{\mathbf L}{\otimes}_S A_{Y_2}),K_{X_1}\overset{\mathbf L}{\otimes}_S L_2) \arrow[d,"\mathrm{triv.}"']
& \\
p'_{1*}\uRHom(L_1\overset{\mathbf L}{\otimes}_S A_{X_2},K_{X_1}\overset{\mathbf L}{\otimes}_S L_2) \arrow[r,"\mathrm{c.b.}"']
& p'_{1*}\uRHom_S(L_1,L_2),
\end{tikzcd}
\]
\[
\begin{tikzcd}[column sep=3.7em,row sep=1.7em]
\uRHom(f_{1!}L_1\overset{\mathbf L}{\otimes}_S A_{X_2},K_{Y_1}\overset{\mathbf L}{\otimes}_S L_2)
  \arrow[r,"\mathrm{c.b.}"] \arrow[d,"\mathrm{c.b.}"']
& \uRHom_S(f_{1!}L_1,L_2) \arrow[ddd,"(4)"] \\
\uRHom(p'_{2!}(L_1\overset{\mathbf L}{\otimes}_S A_{X_2}),K_{Y_1}\overset{\mathbf L}{\otimes}_S L_2) \arrow[d,"\mathrm{dual.}"']
& \\
p'_{2*}\uRHom(L_1\overset{\mathbf L}{\otimes}_S A_{X_2},p_2'^!(K_{Y_1}\overset{\mathbf L}{\otimes}_S L_2)) \arrow[d,"\mathrm{c.b.}"']
& \\
p'_{2*}\uRHom(L_1\overset{\mathbf L}{\otimes}_S A_{X_2},K_{X_1}\overset{\mathbf L}{\otimes}_S L_2) \arrow[r,"\mathrm{c.b.}"']
& p'_{2*}\uRHom_S(L_1,L_2).
\end{tikzcd}
\]
\end{statement}
Isso não é senão uma explicitação de 2.5 no caso particular em que $(u_1,v_1):(L_1,K_{X_1})\to(f_{1!}L_1,K_{Y_1})$ é dado pela identidade de $f_{1!}L_1$ e pelo isomorfismo canônico $f_1^!K_{Y_1}\xrightarrow{\sim}K_{X_1}$, e em que $(u_2,v_2):(A_{X_2},L_2)\to(A_{Y_2},f_{2*}L_2)$ é dado pelo isomorfismo canônico $A_{X_2}\xrightarrow{\sim}f_2^*A_{Y_2}$ e pela identidade de $f_{2*}L_2$.

\begin{statement}{Corolário 3.5}
Suponhamos que $Y_1=Y_2=S$. Seja $u\in\Hom_S(L_1,L_2)$. A imagem direta $f_*u\in\Hom(f_{1!}L_1,f_{2*}L_2)$ (3.3.2) pode ser obtida das duas maneiras equivalentes seguintes:
\end{statement}
\begin{enumerate}[label=\alph*)]
\item Por meio de dualidade trivial, identifica-se $u$ com uma flecha
\[
L_1\longrightarrow p_{1*}p_2^!L_2;
\]
compõe-se com o isomorfismo de mudança de base
\[
p_{1*}p_2^!L_2\xrightarrow{\sim}f_1^!f_{2*}L_2;
\]
$f_*u$ é a adjunta dessa flecha composta.
\item Por meio de adjunção, identifica-se $u$ com uma flecha
\[
p_{2!}p_1^*L_1\longrightarrow L_2;
\]
compõe-se com o isomorfismo de mudança de base
\[
f_2^*f_{1!}L_1\xrightarrow{\sim}p_{2!}p_1^*L_1;
\]
$f_*u$ é a adjunta dessa flecha composta.
\end{enumerate}
A descrição a) (resp. b)) corresponde ao percurso $(f_{1*}(3)\circ(1))^{-1}$ (resp. $f_{2*}(4)\circ(2)$) de 3.4.1.
